\clearpage
\fancyhead[L]{\small DELAMBRE. TABLES ÉCLIPTIQUES.}
\hypertarget{native-608-lower}{}
\begin{center}\small Göttingische gelehrte Anzeigen. No. 45. Pages 441..443. 19 March 1818.\\[3pt]\rule{.72\linewidth}{.3pt}\end{center}

\emph{Tables écliptiques des satellites de jupiter, d'après la théorie de M. le marquis de \textup{\textsc{Laplace}}, et la totalité des observations faites depuis \textup{1662} jusqu'à l'an \textup{1802}; par M. \textup{\textsc{Delambre}}, chevalier de St. Michel etc.} Paris, 1817. Published by the widow Courcier. The introduction, 58 pages. The tables, 32 half-sheets. 4to.

\begingroup\brokenpenalty=10000\relax
The tables for the eclipses of Jupiter's satellites, calculated nearly 30 years ago by Mr. \textsc{Delambre} in accordance with \textsc{Laplace}'s theory and printed in the third edition of \textsc{Lalande}'s Astronomie, were, as is well known, received by astronomers with deserved approval and came into general use. That industrious and indefatigable calculator has since made use of the increased number of observations and the still greater advances in the theory to improve the foundations of his tables afresh. These new tables were the result. They were printed several years ago but have only now been issued. The great geometer under whose direction this laborious work was carried out has already borne honorable testimony to it: `\textsc{Delambre} a exécuté ce travail important avec le plus grand succès; et ses tables qui réprésentent les observations avec l'exactitude des observations mêmes offrent au navigateur un moyen sûr et facile pour avoir sur le champ, par les eclipses des satellites, et surtout par celles du premier, la longitude des lieux où il atterre.' The arrangement of the tables differs little from that of the earlier ones, and the introduction gives very detailed directions for their use. We would only have wished the author to have given a fuller account of the results of his work of improvement, so that one would have been in a position to see clearly and readily what changes the individual underlying constants to be determined from the observations (these constants number, as is well \SourceLeaf{609}known, thirty-two, counting the speed of light) had undergone, and to judge the degree of accuracy to be ascribed to them. From a passage in the introduction one may infer that all these changes were very slight. He declares that, although the errors in the new tables are generally smaller than those in the earlier ones, `ce que nous avons gagné ne vaut peut-être
pas le travail qu'il a couté', quite in the spirit of the view that the value of such a work should be judged solely by its practical utility for the determination of geographical longitude. Only with regard to the coefficient of the equation of light, which Mr. \textsc{Delambre} takes to be $493''2$, does he remark that this determination rests on more than 1000 observations of the first satellite, and that the earlier determination, which is not perceptibly different from it, was based on 500 observations. It is indeed very remarkable that this coefficient gives the aberration constant $20''25$, whereas that constant, when determined from the right ascensions of the Pole Star, comes out somewhat larger. It might, however, be too hasty to go so far as to infer from this a difference between the speed of light from the fixed stars and that from the satellites, or between the speed of light at the circumference of the Earth's orbit and that nearer the Sun. Even so, the matter certainly deserves continued attention.
\par\clearpage\endgroup
